\begin{explanationbox}
	\textbf{\textcolor{CExplanation}{一句话直觉}}

	线面平行 $\Rightarrow$ 过该直线的平面与已知平面的交线平行于该直线。

	\medskip
	\textbf{\textcolor{CExplanation}{核心拆解}}
	\begin{description}[style=nextline, leftmargin=2.8em, labelsep=0.5em, itemsep=0.45em]
		\item[\textbf{要点一（三条件缺一不可）：}] 必须有 $a\parallel\alpha$（线面平行）、$a\subset\beta$（线在面内）、$\alpha\cap\beta=l$（两平面相交），三条共同作用才推出 $a\parallel l$。
		\item[\textbf{要点二（平行关系传递）：}] 该定理实质是将线面平行关系转化为线线平行关系，为证明共面线线平行提供依据。
	\end{description}

	\medskip
	\textbf{\textcolor{CExplanation}{几何本质（平行传递与相交公理）}}

	直线 $a$ 平行于平面 $\alpha$，故 $a$ 与 $\alpha$ 无公共点；$a$ 在平面 $\beta$ 内，$\alpha$ 与 $\beta$ 相交于 $l$，则 $l$ 是 $\beta$ 内唯一与 $\alpha$ 有公共点的直线。若 $a$ 与 $l$ 不平行，则它们在 $\beta$ 内必相交，交点既在 $a$ 上也在 $\alpha$ 上，与 $a\parallel\alpha$ 矛盾。因此 $a\parallel l$。本质是反证法与平行定义的结合。

	\medskip
	\textbf{\textcolor{CExplanation}{代数意义（向量等价刻画）}}

	设平面 $\alpha$ 的法向量为 $\vec{n}$，平面 $\beta$ 的法向量为 $\vec{m}$。由 $a\parallel\alpha$ 得 $a$ 的方向向量 $\vec{v}\perp\vec{n}$；由 $a\subset\beta$ 得 $\vec{v}\perp\vec{m}$。交线 $l=\alpha\cap\beta$ 的方向向量 $\vec{w}$ 同时满足 $\vec{w}\perp\vec{n}$ 和 $\vec{w}\perp\vec{m}$。因为 $\alpha$ 与 $\beta$ 相交，所以 $\vec{n}$ 与 $\vec{m}$ 不共线。在空间中，若一个非零向量同时垂直于两个不共线的向量，则其方向唯一确定。因此 $\vec{v}$ 与 $\vec{w}$ 方向相同，即 $\vec{v}\parallel\vec{w}$，从而 $a\parallel l$。

	\medskip
	\textbf{\textcolor{CExplanation}{考点价值（高频应用场景）}}
	\begin{itemize}[leftmargin=*, itemsep=0.5em, topsep=0.4em]
		\item \textbf{考法一（证明线线平行）：} 已知一条直线与一个平面平行，证明它与某条在另一平面内的直线平行时，可构造过该直线的平面与已知平面相交，得到平行线。
		\item \textbf{考法二（截面问题求交线）：} 在立体几何的截面问题中，若已知截面与某一平面平行，可利用该定理确定截面与各面的交线方向。
	\end{itemize}

	\medskip
	\textbf{\textcolor{CExplanation}{顿悟点}}

	\textbf{构造平面是关键。}当遇到线面平行条件时，优先考虑作一个过该直线的平面与已知平面相交，交线平行于已知直线，从而将空间平行问题转化为平面平行问题。

	\medskip
	\textbf{\textcolor{CExplanation}{使用场景}}
	\begin{itemize}[leftmargin=*, itemsep=0.5em, topsep=0.4em]
		\item \textbf{场景一（已知线面平行求交线）：} 在几何体中出现线面平行条件，需要确定该直线与某一平面的交线时，用过该直线的平面去截。
		\item \textbf{场景二（证明线线平行）：} 若要证明两条直线平行，其中一条与某个平面平行，且另一条是该平面与经过第一条直线的平面的交线，可直接套用定理。
	\end{itemize}
\end{explanationbox}
